\documentclass[letter, 11pt]{article}
\usepackage[utf8]{inputenc}
\usepackage[english]{babel}
\usepackage{amsfonts}
\usepackage{amsmath}
\usepackage[dvips]{graphicx}
\usepackage{url}
\usepackage[top=3cm,bottom=3cm,left=3.5cm,right=3.5cm,footskip=1.5cm,headheight=1.5cm,headsep=.5cm,textheight=3cm]{geometry}
\usepackage{listings}
\usepackage{color}
\usepackage{fancyvrb}
\usepackage{fancyhdr}
\usepackage{enumerate}
\usepackage{amssymb}
\usepackage{multirow}
\usepackage{hyperref}
%% to use the warpfigure capability
\usepackage{wrapfig} 

%%%%%%%%%%%%%%%%%%%%%%
%Estilo del documento%
%%%%%%%%%%%%%%%%%%%%%%
\pagestyle{fancyplain}

%%%%%%%%%%%%%%%%%%%%%%%%%%%%%%%%%%%%%%%%%%%
%Fancyheadings. Top y Bottom del documento%
%%%%%%%%%%%%%%%%%%%%%%%%%%%%%%%%%%%%%%%%%%%
% Recuerde que en este documento la portada del documento no posee
% numeracion, pero de igual manera llamaremos a esa primera pagina la numero
% 1, y la que viene la dos. Esto es para tener una idea de las que
% llamaremos pares e impares
\lhead{} %Parte superior izquierda
\rhead{\bf \it UTFSM, Casa Central} %Parte superior derecha
\lfoot{JSCO} %Parte inferior izquierda.
\cfoot{} %Parte inferior central
\rfoot{\bf \it Page \thepage} %Parte inferior derecha
\renewcommand{\footrulewidth}{0.4pt} %Linea de separacion inferior


\begin{document}
%%%%%%%%%%%%%%%%%%%%%%%%%%%
% Tema de las referencias %
\bibliographystyle{unsrt}
%%%%%%%%%%%%%%%%%%%%%%%%%%

%%%%%%%%%%%%%%%%%%%%%%%%%%
%Definicion de la portada%
%%%%%%%%%%%%%%%%%%%%%%%%%%
\begin{titlepage}
    \begin{center}
	\begin{tabular}{ccc}
	    \includegraphics[width=3cm]{img/utfsm}
	    & 
	    \hspace{-0.2cm}
	    \begin{tabular}{c}
		Universidad T\'ecnica Federico Santa Mar\'ia \\ \hline
		\vspace{0.2cm}
		Departamento de Inform\'atica\\
		\vspace{1.2cm}
	    \end{tabular}
	    \hspace{0.2cm}
	    &
            \includegraphics[width=2cm]{img/di}
	\end{tabular}

	\vspace{1cm}
	%Titulo del Documento
	\begin{tabular}{c}
		\Huge{\sc{Testing, Detection and Possible}}\\\\
		\Huge{\sc{Solutions for the Bufferbloat}}\\\\
		\Huge{\sc{Phenomenon on Local Networks.}}\\\\
		%\Large{\sc{Gemini Observatory}}\\\\
	\end{tabular}

    	\vspace{2cm}
	\Large{\sc{JUAN S. CATALAN OLMOS}}\\
	\vspace{2cm}
	\large{Definici\'on de Tema de Memoria}\\
	\large{para optar al T\'itulo de:}\\
	\large{INGENIERO CIVIL INFORMATICO}\\
	\vspace{2cm}
         	\normalsize{Profesor Gu\'ia: Horst H. von Brand}\\
	%Fecha
    
    \vspace{2cm}
    \normalsize{\sc{\today}}\\
    \normalsize{\sc{Valpara\'iso, Chile}}\\
    \end{center}
\end{titlepage}


\newpage

%section Intro
\section{Introduction}

If a little salt makes food taste better, then a lot must make it taste great, right?. What happens if you apply the same statement to a network domain? It keeps been as good as it was? It improves the performance or makes it worse?.\\

\textit{Lets think of a network as a road system where everyone drives at the maximum speed. When the road gets full, there are only two choices: crash into other cars, or get off the road and wait until things get better. The former isn't as disastrous on a network as it would be in real life: losing packets in the middle of a communication session isn't a big deal. But making a packet wait for a short time is usually better than ``dropping" it and having to wait for a re-transmission.}\cite{bufferarms}\\

At this point, the role of the router becomes important. It has to control the congestion effectively in networks. It is important to remember that the traffic in a network is inherently bursty, so the role of the buffers in the router is to smooth the flow of traffic. Without any buffering, to allocate the bandwidth evenly would be impossible. But there are some problems with current algorithms; they use tail-drop based queue management that has two big drawbacks: 1.- lockout 2.- full queue that impact with a high queue delay.\\

These problems are fixed with the creation of a group of FIFO based queue management mechanisms to support end-to-end congestion control in the internet. That procedure is called Active Queue Management (AQM). With AQM the loss of package and the average queue length is reduced; this impacts in a decreasing end-to-end delay by drooping packages before buffer comes full, using the exponential weighted average queue length as a congestion indicator. For the proper use of AQM, it has to be widely enabled and consistently configured the router.\\

\textit{Today's networks are suffering from unnecessary latency and poor system performance. The culprit is Bufferbloat, the existence of excessively large and frequently full buffers inside the network. Large buffers have been inserted all over the Internet without sufficient thought or testing. They damage or defeat the fundamental congestion-avoidance algorithms of the Internet's most common transport protocol. Long delays from bufferbloat are frequently attributed incorrectly to network congestion, and this misinterpretation of the problem leads to the wrong solutions being proposed.}\cite{GettysNichols}\\

The existence of cheap memory and a misguided desire to avoid packet loss has led to larger and larger buffers being deployed in the hosts, routers, and switches that make up the Internet. It turns out that this is a recipe for bufferbloat. Evidence of bufferbloat has been accumulating over the past decade, but its existence has not yet become a widespread cause for concern.\\




delay-based and loss-based approach

\cite{NicholsJacobsonCQD}
Cheap memory and ``more is better'' mentality have led to the inflation and proliferation of buffers. Second, dinamically varying path characteristics are much more common today and are the norm at the consumer Internet edge.

as each packet is squeezed down in bandwidth, it must stretch out in time since its size stays constant.

This standing queue, resulting from mismatch between the window and pipe size, is the essence of bufferbloat. It creates large delays but no improvement in throughput. It is not a phenomenon treated by queueing or traffic theory, which unfortunately, results in it being almost universally misclassified as congestion.

The bufferbloat problem, making the window match the pipe size, is hard to address. Window sizes are chosen by senders while queues manifest at bottleneck gateways.

%end intro
%section Foundation
\section{The Bufferbloat Foundations}
\subsection{The TCP Protocol \cite{rfc793}\cite{rfc2001}}
%http://www.linktionary.com/c/congestion.html
%http://condor.depaul.edu/jkristof/technotes/congestion.pdf 

The Transport Control Protocol or TCP is main transport protocol of the
internet, providing reliable host to host communication over unreliable
transport media\cite{RFC793}. The advantages of its connectionless design,
flexibility and robustness provides reliable, ordered delivery of a stream of
bytes from a program on one computer to another program on another computer,
but for that, the cost are: the needed of careful design so it provides a good
service under heavy loads. Many modifications has been done since it was
originally defined in 1981.

The root idea of any algorithm related to transport connections must be based
into the ``\textit{packet conservation principle}". This principle claims:

\begin{defn}
A new packet isn't put into the network until an old packet leaves.
\end{defn}


\subsubsection{Slow-Strart Algorithm\cite{rfc879}\cite{rfc2460}}
\indent Old TCPs implementation, would start a connection with the sender injecting multiple segments into the network, up to the window size advertised by the receiver.  This action has no implications if the two host are into the same LAN. While this is OK when the two hosts are on the same LAN; but in the Internet, this schema isn't valid. As is know, between the two end points are routers and getaways, and the flow of packages is not constant, between the two end points some links could be slower than others, and some intermediate buffer's queues could run out of space.\\

The algorithm to starts this ``\textit{clock}'' and to avoid this congestion
 \begin{wrapfigure}{r}{0.5\textwidth}
  \begin{center}
    \includegraphics[width=0.48\textwidth]{img/slowstart}
  \end{center}
  \caption{The Chronology of a Slow-Start.\cite{Jacobson:1988:CAC:52325.52356}}
  \label{slowstart}
\end{wrapfigure}
is called \textit{slow start}. It is the responsible to gradually increase the amount of data that is in transit by observing that the a new package is injected into the network until the acknowledgment of a previous packages as arrives from the other end. In other words, it is used to avoid sending more data than the network is capable of transmit.\\

Slow start adds a variable window to the sender's TCP:  the congestion window, called ``\textit{cwnd}'' to the per-connection state.  When a new connection is established or restarting after a loss connection with a host, the congestion window is initialized to one segment, with the size of two times the maximum segment size (MMS)\footnote{The study of MMS is out of the scope of this paper, but more info could be found in \cite{rfc879} and \cite{rfc2460}} .  Each time an ACK is received, the congestion window is increased by one segment (1 MMS).  The sender can transmit up to the minimum of the congestion window and the advertised window.  The congestion window is flow control imposed by the sender, while the advertised window is flow control imposed by the receiver.  The former is based on the sender's assessment of perceived network congestion; the latter is related to the amount of available buffer space at the receiver for this connection.\\

The sender starts by transmitting one segment and waiting for its ACK. When that ACK is received, the congestion window is incremented from one to two, and two segments can be sent.  When each of those two segments is acknowledged, the congestion window is increased to four as seen in figure \ref{slowstart}. Here, the gray numbered boxes are packages and the white are the corresponding ACK. As each ACK arrives, two packages are generated, one for the ACK package that left the ``\textit{pipe}" and one because an ACK opens the congestion window by one. This provides an exponential growth, it takes time \textbf{$Rlog_2 W$}\cite{Jacobson:1988:CAC:52325.52356}, where R is the round trip time and W is the window size. Although it is not exactly exponential because the receiver may delay its ACKs,typically sending one ACK for every two segments that it receives.\\

At some point the capacity of the internet can be reached, and an intermediate router will start discarding packets.  This tells the sender that its congestion window has gotten too large. Early implementations performed slow start only if the other end was on a different network.  Current implementations always perform slow start.\\


\subsubsection{The Congestion Avoidance Algorithm \cite{rfc2309}\cite{rfc2581}}
\indent In the ``\textit{slow start}" phase, if a when a lost occurs, half of the current window is saved as a Gresham a variable that is
 used to determine whether the slow start or congestion avoidance algorithm is used to control data transmission. After this, the \textit{cwnd} is set again to 1 and start to grown until it reaches the ssthresh again. Now, TCP goes into congestion avoidance mode, where for each ACK increases the cwnd in 1/cwnd. A congestion can occur when data arrives at a router whose output capacity is less than the sum of the inputs.  Congestion avoidance is a way to deal with lost packets.\\
	
This algorithm makes a fundamental assumption: \textit{the packet loss caused by damage is very small (much less than 1\%), therefore the loss of a packet signals congestion somewhere in the network between the source and destination}.  Two are the indication of package lost: a timeout occurring and the receipt of duplicate ACKs.\\

A good congestion avoidance strategy, must have two components: 1.- The endpoints should know when a congestion is about to occur or occurring into the network, and 2.- If a signal that alert the congestion is received, the network's utilization must decrease and increases if the signal isn't received.\\

When a networks is getting congested, the queue lengths will start to increase exponentially (because of slow-start). The system will collapse if the network doesn't throttle back the traffic sources at leas as quick as the queues are growing. The way that the network announces via dropped packets when demand is excessive, but says nothing if a connection is using less than its fair share. This is also another problem because causes underbuffering, also causing that the resources are not fully under full use.\\

The implementation of congestion avoidance is as simple as slow start, but with a slight difference\cite{Jacobson:1988:CAC:52325.52356}. The steps are:
\begin{enumerate}
\item On any timeout, set \textit{cwnd} to half the current window size. This produces a multiplicative decrease.
\item On each ack for new data, increase \textit{cwnd} by 1/\textit{cwnd}. Now cwnd has an additive increment so now the growth becomes linear.
\item When sending, send the minimum of the receiver's advertised window and \textit{cwnd}.
\end{enumerate}

A window of size cwnd packets will generate at most cwnd ACKs in one round trip time. Thus an increment of 1/cwnd per ACK will increase the window by at most one packet in one RTT. In TCP, windows and packet are in bytes so the increment translates to segsize*segsize/cwnd, where segsize is the segment size and cwnd is maintained in bytes.\\

Congestion avoidance and slow start are independent algorithms with different objectives.  But when congestion occurs TCP must slow down  its transmission rate of packets into the network, and then invoke slow start to get things going again. In practice they are implemented together. So, if cwnd is less than or equal to ssthresh, TCP is in slow start; otherwise TCP is performing congestion avoidance. Slow start continues until TCP is halfway to where it was when congestion occurred (since it recorded half of the window size that caused the problem ), and then congestion avoidance takes over.\\


\subsubsection{The Router's Congestion Avodance Complements}
After the ``congestion collapse" and with the growth in the last few decades of
 \begin{wrapfigure}{l}{0.5\textwidth}
  \begin{center}
    \includegraphics[width=0.48\textwidth]{img/overflows}
  \end{center}
  \caption{How Tail-drop management and RED AQM overflows\cite{Dischinger2007CRB}}
  \label{tdredof}
\end{wrapfigure}

Internet, it has become clear that the TCP congestion avoidance mechanisms,
while necessary and powerful, are not enough to provide a fully safe service,
and the control that can be accomplished from the edges of the networks has
proof that has a limit. So, in order to obtain a good service in all
circumstances, some mechanisms in the routers to complement this endpoint
congestion avoidance mechanisms are needed.\\

Two classes of router algorithms related with congestion avoidance control can
be distinguished: ``queue management" and ``scheduling" algorithms. While the
first one manage the length of packet queues by dropping packets when necessary
or appropriate, the second one determine which packet to send next, and are used
to manage the allocation of bandwidth among flows. It is important to notice
that while this two mechanisms are closely related, the performance that they
address are rather different and should be seen as complementary, and not as
replacements for each other.\\



\begin{enumerate}
\item \textbf{Managing the Routers Queue:} As we have seen in previous section, the traditional way to manage router queue
is, after set a maximum length for the queue, accept packages until this length
is reached, and then drop subsequent packages until a packet from the stack as
been transmitted. This technique is called ``tail drop", and has served the
Internet well enough for years, but it has two important drawbacks:\\

\begin{itemize}
\item It allows, in some situations, a single or few flows to monopolize the
queue space, preventing other connections from getting room in the queue.
\item The signaling is produced only when the queue is full, so it allows queue
to maintain almost full status for long periods.
\end{itemize}

If queue is full or almost full, an arriving burst will cause multiple packets to
be dropped, an this behavior can produce a global synchronization of flows
throttling back followed by sustained period of lowered link utilization, which
will impact in a reduce of overall throughput, and if a long flows arrives in
that period, the lock-out of the queue.\\

Besides tail drop, another two techniques can be applied in these situations.
The ``random drop on full" will produce that the router drops a randomly selected
packet from the full queue, which requires an $O(N)$ walk through the queue, when a
new packet arrives. Under the ``drop front on full", the router drops the packet
at the front of the queue. Either of these solve the lock-out problem, but
neither solves the full-queue.\\

\item \textbf{Active Queue Management and Random Early Detection} In the current Internet, dropped packets are used as a critical mechanism to
notify a end node when a congestion is presented. So, if a router is capable to
drop packets before the queue is full, and the end nodes take actions before the
buffers overflow, the full-queue problem is solved. This proactive approach is
know as ``active queue management" (AQM) and allows routers to control when and how
many packets to drop before buffers overflow.\\

For responsive flows, AQM can provide:
\begin{itemize}
\item reduce number of packets dropped in routers
\item provide lower-delay interactive service
\item avoid lock-out behavior
\end{itemize}

One AQM algorithm for routers is called ``Random Early Detection" or RED. The
algorithm drops arriving packets probabilistically, which increases as the
estimated average queue size grows, so it approach is based on the ``recent
past" events.\\

The RED algorithm consists of two main parts:

\begin{itemize}
\item Estimation of the average queue size
\item Packet drop decision
\end{itemize}

RED's particular algorithm for dropping is the culprit in the performance
improvement. 

\end{enumerate}

\subsection{Latency}
When we move an amount of data, like a music file, it takes several minutes, or if we have lucky, several seconds. Smaller the file gets, less time is the duration of the transfer, but there is a limit. No matter how small the file becomes, we are stocked with a minimum time that we can never beat. That is call latency of the device. For an Ethernet network is 0.3\textit{ms}.\\

Maybe we don't notice the effect of this time, that when the amount of data is large enough, this time is too small compared with the time that takes the whole transfer. But, what happens with short-flows, like game streaming? Let's imagine that we want to stream audio over the net. 100\textit{ms} may not sound very much, but it's enough to notice a delay and echo in voice. A better case can be found in \cite{main:ref:3}, where the effects of latency in the transmissions and the buffers can be found.\\

There is no visible impact of varying the latency other than its direct effect of varying the bandwidth-delay product. Congestion can also be caused by deny of service attack that attempts to flood host or routers with large amount of network traffic. \\



%end section Fundation
%section Characterization
\section{Characterization of Bufferbloat}

\subsection{Backbone Routers}

As we already know, all internet routers contains buffers to hold packets during times of congestion. A widely used rule-of-thumb states that each link need a buffer of size $B = \overline{RTT} x C $, where $\overline{RTT}$ is the average round trip time of a flow passing across the link, and C is the data rate of the link. The main characteristic of bufferbloat is the existence of excessively large and frequently full buffers inside the network. Large buffers have been inserted all over the Internet without sufficient thought or testing, so router buffers are the single biggest contributor to uncertainty in the Internet.\\

The rule-of-thumb come from a desire to keep the link as busy as possible so, the throughput of the network is always as big as possible. But, because the way that TCP works, no matter how big the buffer is at the bottleneck link, TCP will cause the buffer to overflow.\\ 

Overbuffering is a bad idea for two reasons:
\begin{enumerate}
\item It complicates the design of high-speed routers, leading to higher power consumption, more board space, and lower density.
\item It increases end-to-end delay in the presence of congestion
\end{enumerate}

As seen in 2.2, large buffers only increases latency, and this only causes conflict with the needs of real time applications.\\

The most important fact of sizing a buffer is to make that sure that while the sender pauses, the router buffer doesn't go empty and force the bottleneck to go idle. Again, the idea is to keep as much throughput as possible so the use of the link is fully utilized. The buffer will avoid to idle if the first packet from the sender shows up at the buffer just as it hits empty. In previous section, we define that after a lost is detected, the cwnd is set to half of is last value, so if we denote as $(W_{max} /2)/C$ the amount of time that packets are sent in congestion phase, and as $B/C$ the time that takes a buffer with size B to drain, the size of a buffer B needed is $B \leq (W_{max} /2)$.\\

Also from \cite{main:ref:1}, we can see that the rule-of-thumb doesn't longer apply to backbone routers, and a better estimator of the size of a buffer with n flows would be no more than $B = (\overline{RTT}xC)/\sqrt{n}$. With the assumption that short-flows plays a very small effect, and that the buffer size is dictated by the number of long flows, this factor will be proof that routers are much longer than they need to be, possible by two order of magnitude.


\subsection{Residential BroadBand Networks}
It is well know that residential networks are often the bottleneck in the last
mile access to the Internet Infrastructure. This could be because the ISP's of
both of the most popular ways to access (DSL and cable networks) to internet
today, use traffic shaping methods and ,as seen in previous sections, deploy
massive queues that can delay packets for several hundred milliseconds.\\

Both shares the asymmetric bandwidths; they downstream bandwidth is higher than
their upstream bandwidth, but in cable networks a single coaxial cable shares
multiple customers, they can concatenate multiple upstream packets into a single
transmission, which result in short bursts at high data rates, so the latency
can heavily fluctuate. This concatenation can produce a jitter time, that under
high network load can be higher than end-to-end jitter over the entire path
under normal load, which can be produce a miss interpretation for some protocols
of incipient congestion and cause to enter into congestion control avoidance too
early.\\

In the other hand, in DSL networks, the maximum data transmission
rate falls with increasing distance from the head, thus in order to boost the
transmission rate, DSL relies on advanced signal processing and error correction
algorithms which can lead to high packet propagation delays .\\




%end section Characterization
%section Test
\section{Experimental Work}
The goal of the experiments outlined in this section is to examine different 
residential and public networks looking for checking the existence of the 
phenomenon under a uncontrolled scenario. This will be done by running a set of 
tests with different tools first to define the characterize and then measure the 
behaviour of the networks with and without any load. More specifically, latency 
under load is tested and analyzed to identify if the latency that occurs is due 
to an excess of buffers or due to some other problem.\\

This section aims to explain the setup that will be used to perform the tests,
the different tools used in each of these tests and what is to be achieved and 
expected from each one of these tests.\\


\subsection{The Test Setup}
The tests were ran under a pseudo controlled environment, using one
physical  machine and a second virtual machine hosted on the first machine. These
systems runs under a regular OS without any modifications. Also, for some tests two other devices will be
added,  one acting as an Iperf Server and a regular Android Tablet that will
be used to add some extra load to the network when the Ethernet cable is used
as medium.

All of these tests will be carried out in a real-world scenario, where no
packet  prioritization is done by the server against our flows, the routes can
vary between each iteration of the same test, and many different flows will
collide with other flows from different sizes and types. Nor there is more
information about how the flows are treated by the queue manager algorithms or
about how they are  configured.

\subsubsection{Hardware Characterization}

\begin{description}

\item [Physical Machine]: \hfill \\
The physical machine runs as host OS, Windows 7 SP1, that works with an Intel(R)
Core(TM) i7-2670QM CPU \@ 2.20GHz with 8GB of available RAM. This machine will 
always be connected through its wireless adapter, a \textit{Broadcom Corp. 
BCM4313 802.11b/g/n Wireless LAN Controller (rev 01)}. The Ethernet controller 
is a \textit{Realtek Semiconductor Co., Ltd. RTL8111/8168 PCI Express Gigabit 
Ethernet controller (rev 06)} adapter, and will be bridged to the virtual 
machine for some tests.

\item[Virtual Machine]: \hfill \\
The virtual machine is hosted using VMWare Player 6.0.1, with 4 processors 
assigned for use plus 4GB of RAM, and with the Ethernet adapter connected only
for certain tests. The OS selected is a Debian based OS called Kali Linux, and
using the kernel release identified as Debian 3.12.6-2kali.

The wireless adapter is a AIR-802 USB adapter with Zydas chipset and a TP-link
8dbi antenna. This adapter is directly connected to the virtual machine and
hooked to the physical machine without the USB extension, this way any extra
signal  loss is avoided. 

\item[Iperf Server]: \hfill \\ 
This is a VPS hosted by Digital Ocean
\footnote{\url{https://digitalocean.com}} with 512MB Ram, 20GB SSD Disk, and
located in New York data center. This machine runs Ubuntu 12.04.3 x64 under
KVM software using as a processor an Intel Hex-Core 3 GHz.

\end{description}

The idea of using an external device to connect the virtual machine to the 
network and not using a bridged configuration provided by software, is mainly 
because with an USB device the machine will take care of all the management and 
administration of the device, avoiding any possibility that the host machine 
modify or manages any flow. Also, by using a second machine to overload the 
uplink, causes to avoid overflow the queue on the machine that is performing 
the tests. With this, it is expected to minimize the possibility that our 
testing machine is causing extra latency, either by the saturation of the 
wireless channel, or by the queue in the traffic control subsystem into the 
kernel. 

For the tests, the Bufferbloat community\cite{bloat} has created a set of best 
practices\cite{tg12} to follow so the results are consistent and repeatable, 
but in this case, computers and routers will not be modified as it will attempt 
to analyze what an every-day-user experience. Those practices will be taken in 
consideration if it is the QOS present in routers and deactivate as recommended.
 


\subsection{Tools Definition}
The tools used for the benchmark were selected by the capability to determinate 
the presence of the Bufferbloat phenomenon in the network of study by the 
analysis of its indicator, the latency or the round trip time\footnote{To 
Subtel, organization responsible for control and supervision in the performance 
of telecommunications in Chile, RTT and Latency correspond to the same thing 
\cite{subtel} }. So, with this as a main consideration, the tool's selected for 
the benchmark will be chosen by the complexity and accuracy to measure and 
determine the RTT into a IP/TCP connection, and how hard is to consistently 
replicate the results under similar contexts. Other tools were selected to help 
to determine to characterize and proof the capability of the network to cause 
Bufferbloat.\\

The tools selected to be used in this tests are the following:

\begin{itemize}
    \item Speedtest test by Ookla
    \item Netalyzr by ICSI
    \item Iperf Tool with Tcptrace/Xplot.org
    \item Page Benchmarker extension for Google Chrome
    \item Smokeping Latency Tool
\end{itemize}

\subsubsection{Speedtest}

Developed by Ookla\footnote{\url{https://www.ookla.com/about}}, this tool is used
for most of the ISP's and many users in Chile to test their broadband's
connections globally. It can be used not only in their website 
\href{http://www.speedtest.net}{\textit{www.speedtest.net}} but also in Android,
iOS or Windows Phone. A command line interface developed in Python can be used 
too for testing internet bandwidth.\\

The server that will be selected to perform every test, either has or not the
fastest ping, is the one hosted by the Pontificia Universidad Cat\'olica de
Valpara\'iso (this host is the default selected most of times by the site also).  

\subsubsection{Netalyzer\cite{netalyzr}}

asdf

\subsubsection{Iperf}
Iperf is a well known and commonly used network testing tool. It can create a
TCP and UDP data streams and measure the bandwidth and the quality of a network
link. It can perform multiple tests like Latency, Jitter or Datagram Loss.\\

Iperf basically tries to send as much information down a connection as quickly
as possible reporting on the throughput achieved. This tool is especially useful
in determining the volume of data that links between two machines can supply.
This two machines define the network, one acting like a server and the second as
the client. For this scenario, the server will be the VPS that only will receive
the Iperf connections (also will be running ssh but without further
interaction). The VM Linux machine will work as or Iperf Client.\\  

As mentioned in Iperf users mailing list \textit{
When one runs TCP tests, there are 2 things that block Iperf from having clear
view of real throughput: buffering on sender's side (TCP/IP stack) and TCP
behaviour itself (acking). What Iperf can measure is the pace with which it
sends data to TCP/IP stack; TCP/IP stack will only accept data from application
when buffers are not full. If the buffer is huge, Iperf will see high
throughput initially, then it will drop. If there's congestion or retrasmission
going on, Iperf will see it as lower throughput}\cite{iperfmaillist}, but the
data generated by Iperf won't be further analyzed because the  idea behind using
this tool is a TCP's packet generator. This means that the packets generated by 
Iperf will captured and analyzed with tcpdump, tcptrace and xplot.org.\\

\subsubsection{Page Benchmarker}

Page Benchmarker is a Google Chrome extension that intent to test page load 
time performance within Chrome. Measures time-to-first-paint, overall page load 
time KB read/written, and several other metrics, and with its capability to 
clear the cache and existing connections between each page load, makes this 
tool one of the main sources results.\\

\subsubsection{Smokeping}
SmokePing is a latency logging and graphing tool that consists of a running 
daemon which organizes the latency measurements and a CGI which presents the 
graphs. With SmokePing give us the ability to measure latency and packet loss 
in the current network, with RRDtool is capable to maintain a long term 
data store and to draw different graphs with the giving up to the minute 
information on the state of each network connection.\\

Smokeping can be configured to perform a wide range of latency measurement 
probes each one directed to an independent target or over a set of targets 
selected for each proof.\\




\subsection{Test Description}
To test the existence of the Bufferbloat phenomenon, five tests are conducted 
as described below. Each test will be repeated under the following contexts:

\begin{enumerate}
\item Twice on the same day in one network to determine if does exist a 
considerable variance in latency for different times in a day.
\item Select different public and private networks with different 
\textit{``speeds''}.
\item Use the Ethernet cable in order to compare the results with those 
previously obtained using Wireless.
\end{enumerate}

Because it is intended to prove the existence of the phenomenon under
circumstances experienced by everyday users, no kernel parameters will be
amended nor changed, and only the  ability to perform QOS on routers that have
this feature is disabled.  Furthermore, the overall question these tests seek
to answer is the following:

\begin{theoremnon}[] 
``The networks that we use every day, have the necessary characteristics to
generate the  Bufferbloat phenomenon whether under low loads and if does
exists, the how  serious are the effects ?''
\end{theoremnon}

\subsubsection{Speed test} 
The idea under this test, is to set a baseline by comparing the speed offered 
by the ISP and the one at the moment of testing. To find the speed provided, 
the  tool used is the Speedtest in the web site of Ookla.

The benefits of using this web site is that not only it will reveal only
the available national uplink and downlink, also it will set the baseline
for the ping. The expected pings, independently of the connection bandwidth,
should be around $\sim12ms$, based on the data presented in \cite{netvtr} and
\cite{netmov} by the two most used ISP in Chile\footnote{Telef\'onica has a
39.2\% market share while VTR owns 38.8\% respectively. \cite{ispshares}}.

As data, it is estimated that the ratio of the average speed as compared to the
rated speed of the service offered for domestic bonds has a variation to 85\%,
while for international links decreases drastically to 35\%\cite{modcompqos}.
Due this factor, none of the tests will be performed on servers located outside
our country, as it is believed that the data will not very representative
regarding actual service.

\subsubsection{Signs of trouble} 
After characterize our network based on the
speed, it is needed to try to define the state of the service based on the
quality in which our network operates, in example: what kind of traffic is
passing through our network, and also get a little more detail on the
average rates of delay and buffers with which the traffic can find along the
path. To collect all this information will be used Netalyzr. 

While laws in Chile assure consumers networks free of traffic shaping
barriers and filters, ISPs apply traffic management
measures\cite{shapemov}\cite{shapevtr} to deliver an optimum experience of the
service's use and thus, achieve an efficient use of the network, thus further to
protect the safety of users while maintaining network stability.

This is why it is expected that some ports or services are partially or
completely blocked, but the most important is that this test will first light on
the existence of the phenomenon under study.

\subsubsection{Collapse test} 
Already having a clear idea of the state of the studied network, the next will
be try to check empirically that the results obtained by the previous test
above described are valid. For this, Iperf will be used to generate as much
traffic as possible for 5 minutes between servers. But the main goal is not to
measure the throughput of the network; instead to capture the packets that
Iperf produces and study them. So before run Iperf, tcpdump will capture this
information. Subsequently, the tcp's trace is taken and extracted with
tcptrace tool and generate the RTT graph.

This exercise will be conducted three times as part of the test, running the
first time as the sole source of network load. For the second and third time,
after 50 seconds after Iperf was started, from the Windows machine will be 
performed an upload to a Dropbox account with a file big enough to the upload 
take more than the time defined to this test. This upload intends to saturate 
the upstream link, leading to an overload of existing buffer (indistinct 
whether the server is national or not, there are several levels where the 
route is common). The upload will be stopped 50 seconds before the time limit.

The time gap established before loading and the end of the test is to let
Iperf frames reach a steady state flow and thus to generate a basis on which
to analyze the time when the traffic is maximum.

The expected RTT are around $\sim12ms$ and $\sim100ms$ without load, for
national and international servers respectively. About the expected value under
load, it is hard to estimate, further that their expected behavior is 
to be stable round a certain value. In case of no stable behavior, the possible 
causes will be analyzed and tracked.

\subsubsection{Load benchmark test}
With the effects of excess buffers already determined, this tests seeks to show
how this phenomenon affects users who surfs through a web browser. This is why
the extension of Google Chrome Page Benchmarker comes in and it will be used to
do 5 iterations of 10 loads to a website and analyze the behavior of these. 

As in the previous experiment, the first will be without any load on the
network. Then, again a file will be uploaded to Dropbox from the second machine
and after 30 seconds from the start of the load, and with the load active, will
proceed to a second iteration. The third will begin the after a minute,
canceling the file upload and after waiting 30 seconds after cancel the upload,
the benchmark will be measured again. The fourth and fifth will be a minute and
a half of a minute since the previous iteration have finished. The site chosen is
\url{http://www.usm.cl}.

As like Jim Gettys showed in his video demonstration\cite{gettysex}, the 
benchmark increases between 10 to 15 times with load against the original times. 
For this test, the expected proportions are lower, this mainly because the 
kernel machine that will run the tests has implemented already some improvements 
to mitigate this problem.

\subsubsection{Smoke the path}
To verify that the effect of Bufferbloat, if it is the case, does not just
happen on a single server or with a single TCP flow, the Smokeping tool will be
configured to perform two types of tests: Fping and echopinghttp. These tests
will be directed to different types of servers (ie: physical and VPS machines)
which have different connection settings and/or type of services (some servers
has dedicated link), and located in both Chile and the United States.

The tool will be left couple of minutes to track and save the state of the
connection under no extra load (only the test machine will be generating
data). After that, a upload of a big file will be performed from a second
machine until it finish or the behavior has become stable and then re-enter to
a phase without load for couple of more minutes and repeat the upload.

With this test, in addition to verifying the validity of the previously captured
data is expected to determine the presence of other factors that can contribute
to the degradation of the quality of service on the network, whether factors
such as quality of service provided by the host, channel/path problems, problems
related to the way handling packets by the router or the machine, and/or any
other that may arise.


%end section Test


%%%%%%%%%%%%%%%%%%%%%%%%%%%
%incluyo referencias
\bibliography{references,rfc}{}
%%%%%%%%%%%%%%%%%%%%%%%%%%%%

\end{document}
